% A signed-kernel theorem for one increasing arm; the product step remains open.
\section{Positive parts of monotone kernels below the positivity threshold}
All sequences are zero-extended. Say that a finite signed sequence has
an initial positive interval if its positive support is $[0,R]$ and
every coefficient after $R$ is nonpositive.

\begin{lemma}[Uniform windows of a signed sequence]
Suppose $p$ has an initial positive interval and $p^+$ is unimodal.
Then $([t]p)^+$ has an initial positive interval and is unimodal.
If $p^+$ is log-concave on its rising side, including its last mode,
then the same is true of $([t]p)^+$.
Moreover, if $p_n=0$ for $n>L$, the positive part of $[t]p$ is
nonincreasing from rank $L$ onward.
\end{lemma}
\begin{proof}
Let $R$ be the last positive input index and write $W=[t]p$.
Through rank $R$, the window agrees with $[t]p^+$, which is
unimodal and has the stated rising-side log-concavity by the
proved uniform-window lemma. For $n>R$,
$W_n-W_{n-1}=p_n-p_{n-t}$.
If $n-t\le R$, the right-hand side is nonpositive; negative indices
are interpreted as zero. If $n-t>R$, every input in the window is
nonpositive, so $W_n\le0$. Thus the positive part is nonincreasing
after $R$ and cannot reappear after vanishing. Joining it to the
unchanged prefix proves unimodality. If the last mode lies at this
joining edge, the next coefficient has only decreased, preserving
the last log-concavity comparison; all earlier rising comparisons
are unchanged. This proves the log-concavity statement too.

For the last assertion, when $n>L$ the same difference is
$-p_{n-t}$. It is nonpositive unless the whole current window is
already nonpositive. Therefore $W^+$ is nonincreasing from $L$.
\end{proof}

\begin{theorem}[All auxiliary peak sizes]
Let $Y$ be increasing, of weight $b$, and let $m\ge1$ be arbitrary.
Set
\[
 D_m(Y)=1+(S_m-1)F_Y,\qquad
 \delta_m(Y)=z^{b+m}D_m(Y)(1/z).
\]
The positive part $\delta_m(Y)^+$ has an initial positive interval,
is unimodal, and is log-concave on its rising side through its last
mode. No condition $m\ge\kappa(Y)-1$ is required for these positive-part
statements. The full kernel is nonnegative precisely at the previously
proved threshold $m\ge\kappa(Y)-1$.
\end{theorem}
\begin{proof}
Induct on the length of $Y$. The empty arm gives $[m]$.
If $m\ge\kappa(Y)-1$, all conclusions follow from the established
nonnegative-kernel theorem, so consider $\kappa(Y)\ge m+2$.
Write $Y=Zt$ and $c=\operatorname{wt}Z$.
The exact threshold recurrence
$\kappa(Y)=\min\{t+1,\kappa(Z)+1\}$ gives
$t>m$ and $\kappa(Z)\ge m+1$.

All ranks of $Z$ through weight $m$ are full. Thus the coefficients
of $D_m(Z)$ at weights $0,1,\ldots,m$ vanish, and
$\delta_m(Z)$ is zero at every index greater than $c-1$.
The reversed append recurrence is
\[
 \delta_m(Y)=[t]\delta_m(Z)-z^{c+m}[t-m].
\]
By induction and the signed-window lemma, the positive part of
$W=[t]\delta_m(Z)$ is unimodal, has an initial positive interval,
is log-concave on its rising side, and is nonincreasing from
rank $c-1$ onward.

The subtracted interval starts at $c+m>c-1$, entirely on this
nonincreasing positive part. Throughout that interval the new
positive part is $(W_n-1)^+$, so it remains nonincreasing.
At its start it can only decrease. The interval ends at $c+t-1$,
whereas $W$ is supported only through $c+t-2$. Thus the upward
jump at the far endpoint is from $-1$ to zero and cannot create a
new positive coefficient. Everything before the subtraction is
unchanged. The positive support, unimodality and rising-side
log-concavity therefore survive. This closes the induction.
\end{proof}

\begin{corollary}[Every signed window of a single reversed arm]
For every increasing arm $Y$ and every $m\ge1$, the positive part
of $([m]-z^m)q_Y$ is unimodal, has an initial positive interval,
and is log-concave on its rising side.
\end{corollary}
\begin{proof}
It differs from $\delta_m(Y)$ only by subtracting the terminal
monomial $z^{b+m}$. The kernel has zero coefficient there, and
all positive coefficients occur earlier. Thus their positive parts
are identical.
\end{proof}

This settles the signed-window assertion for one monotone arm at
every auxiliary size. The analogous assertion for $q_Xq_Y$ is
still unproved; it is not obtained by multiplying two unimodal
positive parts. A finite diagnostic of $25000$ increasing-arm pairs
found no failure for window lengths $1,\ldots,8$, including lengths
below the arm part sizes. That diagnostic motivates the product
question but is not a proof of it.

\begin{theorem}[Additional structure of the positive part]
For every increasing arm and every $m\ge1$, let
$p=\delta_m(Y)^+$. Besides the preceding properties, $p$ is
recursively rising-unimodal in the sense of
\texttt{iterated\_rising\_differences.tex}. It also satisfies
\[
 p_{r+1}p_s\ge p_rp_{s+1}
 \qquad(0\le r\le m-2,\ s\ge r).
\]
The first $m$ coefficients are nondecreasing.
\end{theorem}
\begin{proof}
The first $m$ coefficients are the cumulative coefficients of
$q_Y$, directly from $([m]-z^m)q_Y$; the isolated terminal
correction does not affect them.

For the recursive assertion use the same induction as above.
If the arm weight $b$ is less than $m$, the positive first
differences of the kernel are exactly $q_Y$, followed by zeros:
the cumulative prefix ends at $b$, remains constant to $m-1$,
and is followed by a nonincreasing tail. Thus the single-arm
recursive theorem applies. In the other cases, the proofs of
the admissible kernel theorem and of the signed extension show
that the new positive part agrees with its uniform-window term
until the correction begins, and is nonincreasing thereafter.
Its positive first differences are a terminal truncation of the
positive first differences of that window. The recursive window
lemma and truncation preserve all the asserted orders.

The coefficient comparison is already known in the admissible
case. In the inadmissible step $Y=Zt$, $t>m$, put
$d=\delta_m(Z)$ and let $R$ be the last index with $d_R>0$.
Inductively $d^+$ has the displayed prefix comparison. Its
nonnegative window $W_0=[t]d^+$ has that comparison too, by
causal log-concave convolution. The signed window $W=[t]d$
agrees with $W_0$ through $R$, and $W^+$ is nonincreasing
after $R$. Since the first $m$ input coefficients are positive,
$R\ge m-1$; thus the initial ratios are unchanged and at
least one. Comparisons entirely before $R$ follow from $W_0$.
For the remaining comparisons the later ratio is at most one,
including the step from $R$ to $R+1$; zero endpoints are
handled by cross multiplication. Finally the subtracted interval
begins at $c+m$, after the start of the nonincreasing portion,
and preserves that portion by the already proved splice lemma.
This proves the comparison for the new positive part.
\end{proof}

The unbounded signed-window tail and subtraction-splice inequalities
used here are checked by the Lean kernel in
\texttt{SignedWindowTail.lean}. The composition interpretation,
inductive shape statements, and the complete theorem above remain
written proofs; they are not claimed as a full Lean formalization.
